\section{引言}

\subsection{背景与动机}

近年来，消费级设备上的本地 AI 应用快速增长。用户越来越期望同一台个人设备能够同时提供多个 AI 服务，以满足隐私、响应速度和离线可用性的需求。典型场景包括助手式交互与间歇性突发请求（如图像生成或文档处理）的组合。然而，这类设备仍然高度资源受限：主流笔记本通常只有 8GB 显存，在本文场景中，这不足以让一个常驻 LLM 运行时和一个高 GPU 占用的突发式生成运行时同时常驻。关键难点不只是这两个模型彼此不同，而是它们承担着不同的服务义务：常驻的 LLM 服务在发生竞争时应保持可达，而突发任务可以等待，但一旦被接纳，就应获得排他的加速器访问权。有限专用显存与不对称服务义务之间的错配，促使我们设计一种面向本地 AI 运行时的轻量级调度器。

\subsection{核心挑战}

\textbf{硬性内存约束}。在我们的测试配置中，8GB 显存不足以同时驻留 LLM 运行时（约 6--7GB，包括权重和 KV cache）与图像模型（约 6.7GB），这使传统 GPU 隔离方法难以适用。当一个 GPU 密集型突发任务到来时，常驻任务占用的显存必须迅速释放或降级，否则就会触发 CUDA OOM。

\textbf{不对称的服务需求}。目标场景中的两类工作负载在服务层面行为不同，如表~\ref{tab:task_characteristics} 所示：

\begin{table}[htbp]
\caption{任务特征对比}
\label{tab:task_characteristics}
\centering
\small
\begin{adjustbox}{max width=\columnwidth}
\begin{tabular}{@{}lcccc@{}}
\toprule
任务类型 & 资源需求 & 时延敏感性 & 可排队性 & 语义 \\
 &  &  &  &  \\
\midrule
LLM 对话 & GPU/CPU & 高（流式） & 低 & 常驻 \\
图像生成 & 必须使用 GPU & 中（端到端） & 高 & 一次性 \\
\bottomrule
\end{tabular}
\end{adjustbox}
\end{table}

重点并不是构建一个涵盖所有本地 AI 任务的宽泛分类体系，而是抽取驱动本文调度决策的那种具体不对称性：一种工作负载在竞争下必须保持可达，另一种工作负载则可以排队，但需要独占 GPU 才能高效完成。

\subsection{现有方案的局限}

现有方案主要面向两个不同的设计中心：（1）分布式调度框架（如 Ray、Kubernetes），它们解决的是集群级放置与编排问题，但通常与单节点消费级设备并不匹配；（2）单模型推理优化系统（如 vLLM、FlexGen），它们提升的是单个模型的吞吐或卸载效率，而不能直接处理多个本地服务共享同一块受限 GPU 时的服务级仲裁问题。

与数据中心调度和单模型服务相比，消费级硬件上的本地混合 AI 服务调度问题研究得更少。现有方案并未直接面向 8GB 级设备上“持续可用、可迁移的服务”与“突发式、GPU 独占的作业”之间的竞争，也很少支持在临时降级期间维持可用性。EdgeSched 正是针对这一空白：它强调这种不对称性的控制后果——任务到来时快速释放 GPU 资源，与突发任务结束后恢复高性能服务，其实是两个不同的问题。

\subsection{本文贡献}

本文提出 EdgeSched，一个面向消费级设备混合 AI 工作负载的本地调度框架，其主要贡献包括：

\begin{enumerate}
\item \textbf{面向本地竞争场景的服务角色刻画}。我们使用一种最小化的服务角色不对称性来描述目标场景：常驻的交互式服务在竞争时应保持可达，而突发任务可以排队，但一旦被接纳就需要独占 GPU。我们仅用这种刻画来推导本文研究场景中的调度决策，而非把它当作所有本地 AI 工作负载的通用本体。

\item \textbf{不对称资源迁移策略}。我们表明，对称迁移不仅效率不足，而且在结构上无法同时优化两个方向。下行迁移优先追求快速释放显存以保证内存安全，因此偏向立即中断和激进回收 GPU 状态；上行恢复则优先追求最小化用户可见中断，因此偏向后台预加载和延迟切换。由于这两个目标本身相互牵制，单一的对称策略不可能同时在两条路径上都达到最优。因此，EdgeSched 使用不对称策略并非表面优化，而是由场景约束决定的必要设计。

\item \textbf{面向生命周期的运行时控制}。我们将优先级队列、反压控制、取消语义和崩溃恢复整合在同一个运行时中，以在混合工作负载和失败场景下提供可预测的调度行为。我们的测量显示，语义层面的崩溃恢复约为 44 ms。

\item \textbf{针对实际替代方案的实证验证}。我们在 RTX 5070 Laptop（8GB VRAM）上对系统进行评估，并将其与基于重启的、通用内存管理的以及迁移式的替代方案进行比较，结果表明，这些更简单的设计在我们的部署场景中无法同时提供连续性、低中断恢复以及有界的恢复成本。
\end{enumerate}

EdgeSched 关注的是本地 AI 运行时中一个尚未被充分研究的调度问题；本文以 LLM 服务与图像生成这一组冲突强烈的工作负载为具体实例，讨论在专用显存紧张的条件下如何稳定运行。

\textbf{可复现性}。源代码与实验脚本已包含在补充材料中。